\documentclass[11pt,a4paper]{article}

% --- Core LaTeX Packages ---
\usepackage[utf8]{inputenc}
\usepackage[margin=1in]{geometry}
\usepackage{amsmath,amssymb,amsthm}
\usepackage{xcolor}
\usepackage{listings}
\usepackage{hyperref}
\usepackage{tabularx}
\usepackage{graphicx}

% --- Optional Diagram Packages (include when needed) ---
% \usepackage{forest}
% \usepackage{tikz}

% --- Unified Color Palette ---
\definecolor{primaryblue}{RGB}{24, 75, 140}
\definecolor{codebg}{RGB}{248, 249, 250}
\definecolor{codeframe}{RGB}{210, 215, 222}
\definecolor{codegreen}{RGB}{40, 130, 60}
\definecolor{codeblue}{RGB}{20, 90, 180}
\definecolor{codegray}{RGB}{120, 120, 120}

% --- Hyperref Setup ---
\hypersetup{
    colorlinks=true,
    linkcolor=primaryblue,
    citecolor=primaryblue,
    urlcolor=primaryblue,
    pdftitle={Number Theory},
    pdfauthor={Antigravity Engineering}
}

% --- Theorem & Definition Environments ---
\theoremstyle{definition}
\newtheorem{definition}{Definition}[section]
\newtheorem{theorem}{Theorem}[section]
\newtheorem{lemma}[theorem]{Lemma}
\newtheorem{example}{Example}[section]

% --- Callout Box Environment ---
\newenvironment{calloutbox}[1]{
    \par\vspace{0.3cm}
    \noindent\begin{tabular}{|p{0.96\textwidth}|}
    \hline
    \vspace{0.1cm}
    \textbf{#1}\par\vspace{0.1cm}
}{
    \vspace{0.1cm}
    \\ \hline
    \end{tabular}
    \vspace{0.3cm}\par
}

% --- Modern C++ Code Listing Setup ---
\lstset{
    language=C++,
    basicstyle=\footnotesize\ttfamily,
    keywordstyle=\color{codeblue}\bfseries,
    commentstyle=\color{codegreen}\itshape,
    stringstyle=\color{orange!80!black},
    numberstyle=\tiny\color{codegray},
    numbers=left,
    numbersep=8pt,
    backgroundcolor=\color{codebg},
    frame=single,
    rulecolor=\color{codeframe},
    breaklines=true,
    breakatwhitespace=true,
    tabsize=4,
    showstringspaces=false,
    captionpos=b
}

% --- Title Block ---
\title{\vspace{-1.5cm}\textbf{\Huge Number Theory}\\[0.3cm]
\Large Primes, Divisors, Factorization, GCD, and Euler's Totient Function}
\author{Technical Monograph}
\date{\today}

\begin{document}

\maketitle

\begin{abstract}
\noindent
Number theory is the study of whole numbers and how they divide, multiply, and combine. At its foundation are prime numbers---the indestructible atomic building blocks from which all integers are constructed. This guide starts from absolute zero: what it means for one number to divide another, how to test if a number is prime without checking every possibility, and how to break composite numbers into their unique prime factors. We then explore the Greatest Common Divisor (GCD), the Euclidean Algorithm, divisor-counting formulas, and Euler's Totient function. Each concept includes plain-English explanations, step-by-step worked traces, and modern C++23 implementations designed for solving problems in algorithmic contests.
\end{abstract}

\tableofcontents
\vspace{0.5cm}
\hrule
\vspace{0.5cm}

% ============================================================================
\section{Foundations: Divisibility and Integers}
% ============================================================================

All of number theory is built on integer division: what happens when you divide one whole number by another, and whether there is anything left over.

\subsection{Divisibility and Remainders}
When you divide an integer $a$ by a positive integer $b$, you get a quotient $q$ and a remainder $r$:
\begin{equation}
a = q \cdot b + r, \quad \text{where } 0 \le r < b
\end{equation}

If the remainder $r$ is $0$, we say that \textbf{$b$ divides $a$} (or $b$ is a divisor of $a$, or $a$ is a multiple of $b$). In mathematical notation, we write:
$$ b \mid a \quad \text{(read: ``$b$ divides $a$'')} $$
If $b$ does not divide $a$, we write $b \nmid a$.

\noindent \textbf{Example:} $4 \mid 12$ because $12 = 3 \times 4 + 0$. However, $5 \nmid 12$ because $12 = 2 \times 5 + 2$ (remainder 2).

\subsection{Prime and Composite Numbers}
Every positive integer greater than 1 belongs to one of two categories:

\begin{definition}[Prime Number]
An integer $p > 1$ is \textbf{prime} if its only positive divisors are $1$ and $p$.
\end{definition}

\begin{definition}[Composite Number]
An integer $n > 1$ is \textbf{composite} if it has at least one positive divisor other than $1$ and itself. That is, $n = a \times b$ where $1 < a, b < n$.
\end{definition}

\begin{calloutbox}{What about 0 and 1?}
The number $1$ is \textbf{neither prime nor composite}---it is called a \textbf{unit}. Excluding 1 from being prime is not an arbitrary rule; it is required so that prime factorizations remain unique. (If 1 were prime, you could write $6 = 2 \times 3 = 1 \times 2 \times 3 = 1 \times 1 \times 2 \times 3$, ruining uniqueness). The number $0$ has infinitely many divisors and is neither prime nor composite.
\end{calloutbox}

\noindent The first few prime numbers are:
$$ 2, 3, 5, 7, 11, 13, 17, 19, 23, 29, 31, 37, 41, 43, 47, \dots $$
Notice that \textbf{2 is the only even prime number}. Every other even number is divisible by 2 and therefore composite.

% ============================================================================
\section{Core Concepts: Divisors, Factors, and Factorization}
% ============================================================================

\subsection{The Symmetry of Divisors: The $\sqrt{n}$ Rule}
Suppose you want to find all divisors of $n = 36$. You could test every number from $1$ to $36$. But notice how divisors always appear in pairs:



If $d$ divides $n$, then $n/d$ must also divide $n$. One of these two factors must be less than or equal to $\sqrt{n}$, and the other must be greater than or equal to $\sqrt{n}$. (If both were larger than $\sqrt{n}$, their product would exceed $n$).

\begin{calloutbox}{Where does this come from?}
Think of factors as rectangles. If you have 36 blocks, you can arrange them as a $1 \times 36$ rectangle, a $2 \times 18$ rectangle, and so on. As you make the rectangle wider, it must get shorter. The most "square-like" rectangle you can make is $6 \times 6$. Once you pass the square root (6), the rectangles just repeat but rotated (for example, $9 \times 4$ is just $4 \times 9$ rotated). Therefore, every unique pair is discovered before or exactly at the square root!
\end{calloutbox}

\begin{calloutbox}{The $\sqrt{n}$ Breakthrough}
To find \textbf{all} divisors or to test if a number is prime, you never need to check beyond $\sqrt{n}$. For $n = 1{,}000{,}000$, instead of checking a million numbers, you only check up to $\sqrt{10^6} = 1000$. This changes an impossible problem into an instantaneous calculation.
\end{calloutbox}

\subsection{The Fundamental Theorem of Arithmetic}
Primes are the ``atoms'' of multiplication. Every positive integer greater than 1 can be broken down into a product of primes in exactly one way (ignoring the order of factors).

\begin{theorem}[Fundamental Theorem of Arithmetic]
Every integer $n \ge 2$ can be uniquely expressed in prime-power form:
\begin{equation}
n = p_1^{a_1} \cdot p_2^{a_2} \cdot \dots \cdot p_k^{a_k}
\end{equation}
where $p_1 < p_2 < \dots < p_k$ are distinct prime numbers and $a_i \ge 1$ are positive integer exponents.
\end{theorem}

\noindent \textbf{Example:} For $n = 360$:
$$ 360 = 2^3 \cdot 3^2 \cdot 5^1 $$
There is no other combination of prime powers that multiplies to 360.

\subsection{Counting Divisors: The $\tau(n)$ Function}
How many total divisors does a number have? Instead of listing them all, you can count them directly from the prime factorization.

Any divisor $d$ of $n = p_1^{a_1} p_2^{a_2} \dots p_k^{a_k}$ must take the form:
$$ d = p_1^{b_1} p_2^{b_2} \dots p_k^{b_k}, \quad \text{where } 0 \le b_i \le a_i $$
For each prime $p_i$, you have $(a_i + 1)$ choices for the exponent $b_i$ (from $0$ up to $a_i$). By the Rule of Product, the total number of divisors $\tau(n)$ is:
\begin{equation}
\tau(n) = (a_1 + 1)(a_2 + 1) \dots (a_k + 1) = \prod_{i=1}^k (a_i + 1)
\end{equation}

\begin{calloutbox}{Where does this come from?}
Think of picking a divisor as ordering a meal from a menu. For $360 = 2^3 \cdot 3^2 \cdot 5^1$, you must choose how many $2$s to include (0, 1, 2, or 3), how many $3$s (0, 1, or 2), and how many $5$s (0 or 1). That gives exactly $4 \times 3 \times 2 = 24$ possible combinations. Every unique combination forms a unique divisor.
\end{calloutbox}

\noindent \textbf{Worked Example: Find $\tau(360)$}
\begin{enumerate}
    \item Factorize: $360 = 2^3 \cdot 3^2 \cdot 5^1$.
    \item Add 1 to each exponent: $(3+1), (2+1), (1+1)$.
    \item Multiply: $4 \times 3 \times 2 = \mathbf{24}$ divisors.
\end{enumerate}

\subsection{Sum of Divisors: The $\sigma(n)$ Function}
The sum of all divisors of $n$, denoted $\sigma(n)$, is also directly computable from its prime powers:
\begin{equation}
\sigma(n) = \prod_{i=1}^k \left( 1 + p_i + p_i^2 + \dots + p_i^{a_i} \right) = \prod_{i=1}^k \frac{p_i^{a_i + 1} - 1}{p_i - 1}
\end{equation}

\begin{calloutbox}{Where does this come from?}
If you expand the product $(1 + p_1 + p_1^2)(1 + p_2)$, you get $1 + p_2 + p_1 + p_1 p_2 + p_1^2 + p_1^2 p_2$. This algebraic expansion perfectly generates every possible divisor exactly once! So, the sum of all divisors is just the product of the sums of prime powers. Since $1 + p + p^2 + \dots + p^a$ is a geometric series, we can use the formula $(p^{a+1} - 1)/(p - 1)$ to compute it instantly.
\end{calloutbox}

\noindent \textbf{Worked Example: Find $\sigma(360)$}
\begin{enumerate}
    \item Factorize: $360 = 2^3 \cdot 3^2 \cdot 5^1$.
    \item Evaluate for $p=2$: $(2^4 - 1)/(2 - 1) = 15 / 1 = 15$.
    \item Evaluate for $p=3$: $(3^3 - 1)/(3 - 1) = 26 / 2 = 13$.
    \item Evaluate for $p=5$: $(5^2 - 1)/(5 - 1) = 24 / 4 = 6$.
    \item Multiply: $15 \times 13 \times 6 = \mathbf{1170}$.
\end{enumerate}

\subsection{Perfect Numbers: When $\sigma(n) = 2n$}
A \textbf{perfect number} is a positive integer equal to the sum of all its \textbf{proper} divisors (every divisor except the number itself). Equivalently, $n$ is perfect when $\sigma(n) = 2n$ --- because the sum of \emph{all} divisors equals $n$ (the number itself) plus the sum of the proper divisors.

\noindent \textbf{Example: Is 28 a perfect number?}
\begin{enumerate}
    \item Factorize: $28 = 2^2 \cdot 7$.
    \item Find all divisors: $1, 2, 4, 7, 14, 28$. Proper divisors (excluding 28): $\{1, 2, 4, 7, 14\}$.
    \item Sum of proper divisors: $1 + 2 + 4 + 7 + 14 = \mathbf{28}$. Yes, 28 is perfect.
    \item Verification using $\sigma$: $\sigma(28) = (1+2+4)(1+7) = 7 \times 8 = 56 = 2 \times 28$. Confirmed.
\end{enumerate}

\begin{calloutbox}{The Perfect Number Test in Code}
To test if $n$ is perfect, compute the sum of divisors using the $O(\sqrt{n})$ pair-collection method (stopping before adding $n$ itself), then check if the sum equals $n$. By convention, 1 is not considered perfect because it has no proper divisors (the sum is 0).
\end{calloutbox}

The first four perfect numbers are $6, 28, 496, 8128$. All known even perfect numbers come from the formula $2^{p-1}(2^p - 1)$ when $2^p - 1$ is prime (a \textbf{Mersenne prime}). Whether any odd perfect numbers exist is one of the oldest unsolved problems in mathematics.

% ============================================================================
\section{Greatest Common Divisor (GCD) and LCM}
% ============================================================================

\subsection{Greatest Common Divisor (GCD)}
The \textbf{Greatest Common Divisor} of two integers $a$ and $b$ (written $\gcd(a, b)$) is the largest positive integer that divides both $a$ and $b$ with zero remainder.

\noindent \textbf{Worked Example: Find $\gcd(24, 36)$}
\begin{enumerate}
    \item Divisors of 24: $1, 2, 3, 4, 6, 8, 12, 24$
    \item Divisors of 36: $1, 2, 3, 4, 6, 9, 12, 18, 36$
    \item Common divisors: $1, 2, 3, 4, 6, 12$. The largest is $\mathbf{12}$.
\end{enumerate}

\begin{definition}[Coprime / Relatively Prime]
Two integers $a$ and $b$ are \textbf{coprime} (or relatively prime) if their only common divisor is $1$:
\begin{equation}
\gcd(a, b) = 1
\end{equation}
\end{definition}

\noindent \textbf{Example:} $\gcd(8, 15) = 1$. Even though neither 8 nor 15 is prime, they share no common prime factors ($8 = 2^3$, $15 = 3 \times 5$). Thus, they are coprime.

\subsection{The Euclidean Algorithm}
The naive way to find $\gcd(a, b)$ is to factorize both numbers, which takes $O(\sqrt{n})$ time. The \textbf{Euclidean Algorithm} finds the GCD in logarithmic time $O(\log \min(a, b))$ using a simple observation:
\begin{equation}
\gcd(a, b) = \gcd(b, a \bmod b), \quad \text{with base case } \gcd(a, 0) = a
\end{equation}

\textbf{Why does this work?} If a number $d$ divides both $a$ and $b$, then $d$ must also divide their difference $a - q \cdot b$, which is the remainder $a \bmod b$. The set of common divisors does not change when replacing $a$ with $a \bmod b$.

\begin{calloutbox}{Where does this come from?}
Imagine a $252 \times 105$ rectangle. You want to tile it perfectly with the largest possible square tiles. First, you place two $105 \times 105$ squares. The uncovered part is a $105 \times 42$ rectangle. Now tile that with $42 \times 42$ squares. You fit two, leaving a $42 \times 21$ rectangle. Finally, you fit two $21 \times 21$ squares perfectly. The side length of the final square that perfectly tiles the space (21) is the GCD! This geometric process perfectly mirrors $\gcd(252, 105) \to \gcd(105, 42) \to \gcd(42, 21) \to \gcd(21, 0)$.
\end{calloutbox}

\noindent \textbf{Worked Example: Euclidean Algorithm Trace for $\gcd(252, 105)$}

\begin{enumerate}
    \item $a = 252, b = 105$. $252 \bmod 105 = 42$. Next step: $\gcd(105, 42)$.
    \item $a = 105, b = 42$. $105 \bmod 42 = 21$. Next step: $\gcd(42, 21)$.
    \item $a = 42, b = 21$. $42 \bmod 21 = 0$. Next step: $\gcd(21, 0)$.
    \item Since $b = 0$, the GCD is $\mathbf{21}$.
\end{enumerate}

\subsection{Least Common Multiple (LCM)}
The \textbf{Least Common Multiple} of $a$ and $b$ (written $\operatorname{lcm}(a, b)$) is the smallest positive integer that is a multiple of both $a$ and $b$.

There is a fundamental connection between GCD and LCM:
\begin{equation}
\gcd(a, b) \times \operatorname{lcm}(a, b) = a \times b
\end{equation}
Rearranging gives the standard formula used in code:
\begin{equation}
\operatorname{lcm}(a, b) = \frac{a \times b}{\gcd(a, b)} = \frac{a}{\gcd(a, b)} \times b
\end{equation}
In programming, always divide by the GCD \emph{before} multiplying to prevent integer overflow.

\begin{calloutbox}{The GCD×LCM Identity: The Most Useful Shortcut}
For any two positive integers $a$ and $b$:
$$ \gcd(a, b) \times \operatorname{lcm}(a, b) = a \times b $$
This means if you know any two of the three values (GCD, LCM, or the product), you can compute the third instantly. In code: always compute $\operatorname{lcm}(a, b) = (a / \gcd(a, b)) \times b$ --- divide \emph{first}, then multiply, to avoid overflow.
\end{calloutbox}

\begin{calloutbox}{Pitfall: Integer Overflow in LCM}
Even if $a$ and $b$ fit in 32-bit integers, their product $a \times b$ might exceed $2^{31}-1$. Computing $(a \times b) / \gcd(a, b)$ can overflow during the multiplication step. By calculating $(a / \gcd(a, b)) \times b$, you keep the intermediate values as small as possible. However, the final LCM can still exceed 32 bits, so always use 64-bit integers (\texttt{int64\_t} or \texttt{long long}) when computing LCM.
\end{calloutbox}

\noindent \textbf{Worked Example: Compute $\operatorname{lcm}(252, 105)$}
\begin{enumerate}
    \item We know $\gcd(252, 105) = 21$.
    \item Divide first: $252 / 21 = 12$.
    \item Multiply: $12 \times 105 = \mathbf{1260}$.
    \item Therefore, $1260$ is the smallest number divisible by both $252$ and $105$.
\end{enumerate}

\begin{calloutbox}{Real-World Analogy: GCD as Simplification}
Imagine you have the fraction $\frac{12}{18}$ and you want to simplify it. You divide both numerator and denominator by their GCD (which is 6): $\frac{12}{18} = \frac{12/6}{18/6} = \frac{2}{3}$. The GCD is exactly the largest number by which you can shrink both sides without creating fractions. In code, you do this any time you reduce a ratio or check if one thing divides another cleanly.
\end{calloutbox}

% ============================================================================
\section{Euler's Totient Function $\phi(n)$}
% ============================================================================

\subsection{What Does the Totient Function Count?}
For any positive integer $n$, \textbf{Euler's Totient Function}, written $\phi(n)$ (Greek letter Phi), counts how many integers in the range $[1, n]$ are \textbf{coprime to $n$}:
\begin{equation}
\phi(n) = \big| \{ k \in \{1, 2, \dots, n\} \mid \gcd(k, n) = 1 \} \big|
\end{equation}

\noindent \textbf{Small Example:} What is $\phi(9)$?
The numbers from 1 to 9 are: $\{1, 2, 3, 4, 5, 6, 7, 8, 9\}$.
\begin{itemize}
    \item Numbers sharing a factor with 9 (multiples of 3): $\{3, 6, 9\}$.
    \item Numbers coprime to 9: $\{1, 2, 4, 5, 7, 8\}$.
\end{itemize}
There are 6 such numbers, so $\phi(9) = \mathbf{6}$.

\subsection{The Formula for $\phi(n)$}
Calculating $\phi(n)$ by checking $\gcd(k, n)$ for every $k$ takes $O(n \log n)$ time. Instead, Euler proved an exact product formula based on prime factorization:

\begin{theorem}[Euler's Totient Formula]
For an integer $n$ with prime factorization $n = p_1^{a_1} p_2^{a_2} \dots p_k^{a_k}$:
\begin{equation}
\phi(n) = n \cdot \prod_{p \mid n} \left( 1 - \frac{1}{p} \right) = n \cdot \left(1 - \frac{1}{p_1}\right) \left(1 - \frac{1}{p_2}\right) \dots \left(1 - \frac{1}{p_k}\right)
\end{equation}
\end{theorem}

\begin{calloutbox}{Where does this come from?}
Think in terms of probabilities or fractions:
\begin{itemize}
    \item A fraction $1/p$ of all numbers up to $n$ are multiples of $p$.
    \item Therefore, a fraction $(1 - 1/p)$ are \emph{not} divisible by $p$.
    \item Multiplying these fractions across all distinct prime factors removes all numbers that share any prime factor with $n$.
\end{itemize}
\end{calloutbox}

\noindent \textbf{Worked Example: Computing $\phi(72)$ from Prime Factorization}
\begin{enumerate}
    \item Factorize: $72 = 2^3 \cdot 3^2$. Distinct primes: $\{2, 3\}$.
    \item Apply the product formula:
    $$ \phi(72) = 72 \cdot \left(1 - \frac{1}{2}\right) \cdot \left(1 - \frac{1}{3}\right) = 72 \cdot \frac{1}{2} \cdot \frac{2}{3} = 72 \cdot \frac{1}{3} = \mathbf{24}. $$
    \item Verification: Check a few values. Numbers from 1 to 72 that share a factor with 72 are multiples of 2 ($36$ of them) plus multiples of 3 not already counted ($12$ of them), totalling $48$ excluded. So $72 - 48 = 24$ remain coprime to 72.
\end{enumerate}

\subsection{Special Cases and Key Properties}
\begin{enumerate}
    \item \textbf{If $p$ is prime:} All numbers from $1$ to $p-1$ are coprime to $p$:
    $$ \phi(p) = p - 1 $$
    \item \textbf{Prime powers:} Only multiples of $p$ share a factor with $p^k$. There are $p^k / p = p^{k-1}$ multiples:
    $$ \phi(p^k) = p^k - p^{k-1} = p^k \left( 1 - \frac{1}{p} \right) $$
    \item \textbf{Multiplicative Property:} If $\gcd(a, b) = 1$, then:
    $$ \phi(a \cdot b) = \phi(a) \cdot \phi(b) $$
    \item \textbf{Sum of Totients of Divisors:} Summing $\phi(d)$ over all divisors of $n$ returns $n$:
    $$ \sum_{d \mid n} \phi(d) = n $$
\end{enumerate}

% ============================================================================
\section{Key Algorithmic Patterns and Techniques}
% ============================================================================

In algorithmic problem-solving, numbers can reach $10^9$ or $10^{18}$. A naive algorithm that counts one-by-one will immediately cause Time Limit Exceeded (TLE). Here are the foundational algorithmic patterns of number theory.

\begin{calloutbox}{The One Idea Behind Every Efficient Number Theory Algorithm}
If $n$ has any factor other than 1 and itself, it must have at least one factor that is $\le \sqrt{n}$. This single observation is the engine behind trial division, the Sieve of Eratosthenes, prime factorization, and the segmented sieve. Every efficient algorithm in this section exploits it. When you see a loop that runs up to $\sqrt{n}$, this is why.
\end{calloutbox}

\subsection{Why Naive Approaches Fail}
\begin{itemize}
    \item \textbf{Checking every number up to $n$:} Testing if $n = 10^9 + 7$ is prime by dividing by all numbers up to $n$ requires $10^9$ operations ($\approx 1$ second in C++). Doing this for 100 queries takes minutes.
    \item \textbf{Generating primes individually:} Calling an $O(\sqrt{n})$ primality test on every number up to $N = 10^7$ takes $O(N \sqrt{N}) \approx 10^{10}$ operations.
    \item \textbf{Factoring repeatedly:} If a problem asks you to factorize $10^5$ queries, running $O(\sqrt{n})$ on each query takes $10^5 \times 10^4 = 10^9$ operations.
\end{itemize}

\subsection{Pattern 1: The Sieve of Eratosthenes (Batch Prime Finding)}
Instead of testing each number for primality, start with an array of numbers from $2$ to $N$ marked as prime. Find the first prime (2), and cross out all of its multiples ($4, 6, 8, \dots$). Move to the next uncrossed number (3), and cross out its multiples ($6, 9, 12, \dots$). Repeat until $\sqrt{N}$.

\vspace{0.2cm}
\noindent \textbf{Key Optimization:} When crossing out multiples of $p$, start at $p^2$ rather than $2p$, because all smaller multiples ($2p, 3p, \dots$) have already been crossed out by smaller prime factors!

\begin{calloutbox}{Pitfall: Inner Loop Starting at $i \times i$}
A common mistake is writing the inner loop as \texttt{for (int j = i * 2; ...)}. This does unnecessary work. Starting at \texttt{i * i} is correct and drastically faster, but beware of integer overflow if \texttt{i} is near $10^5$. Always use \texttt{int64\_t} for the loop variables if $N$ is large!
\end{calloutbox}

\vspace{0.2cm}
\noindent \textbf{Complexity:} $O(N \log \log N)$ time, which is practically linear. For $N = 10^7$, this runs in under 0.1 seconds in C++.

\noindent \textbf{Worked Trace: Sieve of Eratosthenes for $N = 30$}
{\footnotesize

}
We stop because the next prime is $7 > \sqrt{30}$. Remaining primes: $\mathbf{2, 3, 5, 7, 11, 13, 17, 19, 23, 29}$.

\subsection{Pattern 2: The Smallest Prime Factor (SPF) / Linear Sieve}
When you have multiple factorization queries ($Q \le 10^6$), you can precompute the \textbf{Smallest Prime Factor} (SPF) for every number up to $N$.
\begin{itemize}
    \item \texttt{spf[x]} stores the smallest prime that divides $x$.
    \item To factorize any number $x$, repeatedly divide by \texttt{spf[x]}:
    $$ x \leftarrow x / \texttt{spf}[x] $$
    \item Each division takes $O(1)$ time, so factorization takes $O(\log x)$ time!
\end{itemize}

\noindent \textbf{Worked Trace: Factorize 60 using the SPF array (with $N = 60$)}


After running the SPF sieve, the relevant entries are: \texttt{spf[60]=2}, \texttt{spf[30]=2}, \texttt{spf[15]=3}, \texttt{spf[5]=5}.
\begin{enumerate}
    \item $x = 60$. \texttt{spf[60] = 2}. Divide: $60 / 2 = 30$, $30 / 2 = 15$. Factor $2^2$ extracted. $x = 15$.
    \item $x = 15$. \texttt{spf[15] = 3}. Divide: $15 / 3 = 5$. Factor $3^1$ extracted. $x = 5$.
    \item $x = 5$. \texttt{spf[5] = 5}. Divide: $5 / 5 = 1$. Factor $5^1$ extracted. $x = 1$.
    \item Result: $60 = 2^2 \cdot 3 \cdot 5$. Total steps: 4 divisions.
\end{enumerate}
Compare this to trial division on 60, which would test divisors $2, 3, 5, 7 (> \sqrt{60})$ --- 4 candidates. For larger $x$ (say $x = 2^{20}$), SPF takes only 20 lookups versus trial division's $\sqrt{2^{20}} \approx 1024$ attempts.

\subsection{Pattern 3: The Segmented Sieve (Primes in Range $[L, R]$)}
What if you need to find primes in the range $[L, R]$ where $L, R \le 10^{12}$, but the range length $R - L \le 10^6$?
\begin{itemize}
    \item You cannot allocate an array of size $10^{12}$ (Memory Limit Exceeded).
    \item But any composite number in $[L, R]$ must have a prime factor $\le \sqrt{R} \le 10^6$.
    \item \textbf{Technique:} Sieve the small primes up to $\sqrt{R} = 10^6$ first. Then allocate an array of size $(R - L + 1)$ and use the small primes to cross out multiples inside $[L, R]$.
\end{itemize}

\noindent \textbf{Worked Trace: Segmented Sieve Concept for Primes in {$[20, 35]$}}
\begin{enumerate}
    \item Find small primes up to $\lfloor \sqrt{35} \rfloor = 5$. Small primes are $\{2, 3, 5\}$.
    \item Create a boolean array of length $R - L + 1 = 35 - 20 + 1 = 16$ representing indices $0$ to $15$ (values $20$ to $35$), initially all marked as prime.
    \item Strike multiples of $p = 2$: First multiple $\ge 20$ is $20$. Multiples: $20, 22, 24, 26, 28, 30, 32, 34$.
    \item Strike multiples of $p = 3$: First multiple $\ge 20$ is $21$. Multiples: $21, 24, 27, 30, 33$.
    \item Strike multiples of $p = 5$: First multiple $\ge 20$ is $20$. Multiples: $20, 25, 30, 35$.
    \item Remaining uncrossed values in $[20, 35]$ are: $\mathbf{23, 29, 31}$.
\end{enumerate}

% ============================================================================
\section{Worked Examples with Step-by-Step Traces}
% ============================================================================

\begin{example}[Finding All Divisors of $n = 36$]
Find all positive divisors of $36$ using the square root pairing method.
\begin{enumerate}
    \item Compute limit: $\lfloor \sqrt{36} \rfloor = 6$.
    \item Iterate $i$ from $1$ to $6$:
    \begin{itemize}
        \item $i = 1$: $36 \bmod 1 = 0 \implies$ pair is $(1, 36/1) = (\mathbf{1}, \mathbf{36})$.
        \item $i = 2$: $36 \bmod 2 = 0 \implies$ pair is $(2, 36/2) = (\mathbf{2}, \mathbf{18})$.
        \item $i = 3$: $36 \bmod 3 = 0 \implies$ pair is $(3, 36/3) = (\mathbf{3}, \mathbf{12})$.
        \item $i = 4$: $36 \bmod 4 = 0 \implies$ pair is $(4, 36/4) = (\mathbf{4}, \mathbf{9})$.
        \item $i = 5$: $36 \bmod 5 = 1 \implies$ not a divisor.
        \item $i = 6$: $36 \bmod 6 = 0 \implies$ pair is $(6, 36/6) = (6, 6) \implies$ add $\mathbf{6}$ once.
    \end{itemize}
    \item Sorted list of divisors: $[1, 2, 3, 4, 6, 9, 12, 18, 36]$. Total = 9 divisors.
\end{enumerate}
\end{example}

\begin{example}[Prime Factorization and Divisor Analysis for $n = 360$]
Factorize $360$ into prime powers, then compute its total number and sum of divisors.
\begin{enumerate}
    \item Divide out 2: $360 / 2 = 180 \implies 180 / 2 = 90 \implies 90 / 2 = 45$. Count of 2s = 3.
    \item Divide out 3: $45 / 3 = 15 \implies 15 / 3 = 5$. Count of 3s = 2.
    \item Divide out 5: $5 / 5 = 1$. Count of 5s = 1.
    \item Prime factorization: $360 = 2^3 \cdot 3^2 \cdot 5^1$.
    \item Count of divisors $\tau(360) = (3 + 1)(2 + 1)(1 + 1) = 4 \times 3 \times 2 = \mathbf{24}$.
    \item Sum of divisors $\sigma(360) = (1+2+4+8)(1+3+9)(1+5) = 15 \times 13 \times 6 = \mathbf{1170}$.
\end{enumerate}
\end{example}

\begin{example}[Euclidean Algorithm Trace: $\gcd(252, 105)$ and $\operatorname{lcm}(252, 105)$]
Compute the GCD and LCM of $252$ and $105$ using Euclidean modulo reduction.
\begin{enumerate}
    \item $\gcd(252, 105)$: $252 \bmod 105 = 42 \implies \gcd(105, 42)$.
    \item $\gcd(105, 42)$: $105 \bmod 42 = 21 \implies \gcd(42, 21)$.
    \item $\gcd(42, 21)$: $42 \bmod 21 = 0 \implies \gcd(21, 0) = \mathbf{21}$.
    \item Compute LCM using $\operatorname{lcm}(a, b) = (a / \gcd(a, b)) \times b$:
    $$ \operatorname{lcm}(252, 105) = \left( \frac{252}{21} \right) \times 105 = 12 \times 105 = \mathbf{1260}. $$
\end{enumerate}
\end{example}

\begin{example}[Computing Euler's Totient $\phi(36)$]
Find how many integers in $[1, 36]$ are coprime to $36$.
\begin{enumerate}
    \item Find distinct prime factors of $36$: $36 = 2^2 \cdot 3^2$. The distinct primes are $p = 2$ and $p = 3$.
    \item Apply Euler's product formula:
    $$ \phi(36) = 36 \cdot \left(1 - \frac{1}{2}\right) \left(1 - \frac{1}{3}\right) = 36 \cdot \frac{1}{2} \cdot \frac{2}{3} = 36 \cdot \frac{1}{3} = \mathbf{12}. $$
    \item Verification by listing: The 12 coprime integers are $1, 5, 7, 11, 13, 17, 19, 23, 25, 29, 31, 35$.
\end{enumerate}
\end{example}

\begin{example}[Segmented Sieve Concept: Primes in {$[20, 35]$}]
Find all prime numbers between $L = 20$ and $R = 35$.
\begin{enumerate}
    \item Find small primes up to $\lfloor \sqrt{35} \rfloor = 5$. Small primes are $\{2, 3, 5\}$.
    \item Create a boolean array of length $R - L + 1 = 35 - 20 + 1 = 16$ representing indices $0$ to $15$ (values $20$ to $35$), initially all marked as prime.
    \item Strike multiples of $p = 2$: First multiple $\ge 20$ is $20$. Multiples: $20, 22, 24, 26, 28, 30, 32, 34$.
    \item Strike multiples of $p = 3$: First multiple $\ge 20$ is $21$. Multiples: $21, 24, 27, 30, 33$.
    \item Strike multiples of $p = 5$: First multiple $\ge 20$ is $20$. Multiples: $20, 25, 30, 35$.
    \item Remaining uncrossed values in $[20, 35]$ are: $\mathbf{23, 29, 31}$.
\end{enumerate}
\end{example}

% ============================================================================
\section{C++ Implementations for Algorithmic Challenges}
% ============================================================================

All implementations use modern \textbf{C++23} standards, avoiding obsolete practices like raw pointers or global variables.

\subsection{What is Trial Division?}
Trial division is the direct process of testing whether a candidate factor divides $n$. Instead of checking all numbers up to $n$, we check only up to $\sqrt{n}$. Furthermore, after checking 2 and 3, any remaining prime must be of the form $6k \pm 1$ ($5, 7, 11, 13, 17, 19, \dots$), allowing us to skip all other even numbers and multiples of 3.

\begin{lstlisting}[caption={$O(\sqrt{n})$ Primality Testing with 6k+-1 Optimization}]
#include <cstdint>

[[nodiscard]] constexpr bool isPrime(int64_t n) noexcept {
    // 1. Handle base edge cases
    if (n <= 1) return false;
    if (n <= 3) return true; // 2 and 3 are prime
    if (n % 2 == 0 || n % 3 == 0) return false; // Exclude multiples of 2 and 3

    // 2. Test candidate factors of the form 6k - 1 and 6k + 1 up to sqrt(n)
    for (int64_t i = 5; i * i <= n; i += 6) {
        if (n % i == 0 || n % (i + 2) == 0) {
            return false;
        }
    }
    return true;
}
\end{lstlisting}

\subsection{What is the Divisor Extraction Algorithm?}
By iterating from $1$ up to $\sqrt{n}$, each time we find a divisor $i$, we automatically obtain its matching partner $n / i$. Collecting both and sorting yields all divisors.

\begin{lstlisting}[caption={Extracting All Divisors in $O(\sqrt{n})$ Sorted Order}]
#include <vector>
#include <algorithm>
#include <cstdint>

[[nodiscard]] std::vector<int64_t> getDivisors(int64_t n) {
    std::vector<int64_t> divisors;
    if (n <= 0) return divisors;

    // 1. Iterate up to sqrt(n)
    for (int64_t i = 1; i * i <= n; ++i) {
        if (n % i == 0) {
            divisors.push_back(i);
            if (i * i != n) { // Avoid adding the square root twice for perfect squares
                divisors.push_back(n / i);
            }
        }
    }

    // 2. Sort to return in ascending order (using C++23 ranges)
    std::ranges::sort(divisors);
    return divisors;
}
\end{lstlisting}

\subsection{What is Prime Factorization?}
To factorize $n$, divide out 2 until $n$ is odd, then check odd candidates $3, 5, 7, \dots$ up to $\sqrt{n}$. Whenever a candidate divides $n$, divide it out completely and count the exponent. If $n > 1$ remains after the loop, that remaining value is itself a prime.

\begin{lstlisting}[caption={Prime Factorization in $O(\sqrt{n})$ Time}]
#include <vector>
#include <utility>
#include <cstdint>

[[nodiscard]] std::vector<std::pair<int64_t, int>> primeFactorize(int64_t n) {
    std::vector<std::pair<int64_t, int>> factors;

    // 1. Extract power of 2
    if (n % 2 == 0) {
        int count = 0;
        while (n % 2 == 0) {
            ++count;
            n /= 2;
        }
        factors.emplace_back(2, count);
    }

    // 2. Extract odd prime factors up to sqrt(n)
    for (int64_t d = 3; d * d <= n; d += 2) {
        if (n % d == 0) {
            int count = 0;
            while (n % d == 0) {
                ++count;
                n /= d;
            }
            factors.emplace_back(d, count);
        }
    }

    // 3. If n > 1 remains, the remaining number is prime
    if (n > 1) {
        factors.emplace_back(n, 1);
    }

    return factors;
}
\end{lstlisting}

\subsection{How Does the Euclidean Algorithm Work in Code?}
The Euclidean algorithm replaces $(a, b)$ with $(b, a \bmod b)$ until $b = 0$. In C++17 onward, \texttt{std::gcd} and \texttt{std::lcm} are included in \texttt{<numeric>}. Here is the production implementation showing the internal logic.

\begin{lstlisting}[caption={Iterative Euclidean Algorithm for GCD and LCM}]
#include <cstdint>

[[nodiscard]] constexpr int64_t gcd(int64_t a, int64_t b) noexcept {
    while (b != 0) {
        int64_t rem = a % b;
        a = b;
        b = rem;
    }
    return a < 0 ? -a : a; // Return non-negative GCD
}

[[nodiscard]] constexpr int64_t lcm(int64_t a, int64_t b) noexcept {
    if (a == 0 || b == 0) return 0;
    // Divide before multiplying to prevent 64-bit integer overflow
    return (a / gcd(a, b)) * b;
}
\end{lstlisting}

\subsection{What is the Sieve of Eratosthenes?}
The Sieve creates a boolean table representing numbers $0 \dots N$. Starting from 2, whenever an uncrossed prime $p$ is encountered, its multiples $p^2, p^2 + p, p^2 + 2p, \dots$ are crossed out.

\begin{lstlisting}[caption={Sieve of Eratosthenes Generating Primes up to N}]
#include <vector>
#include <cstdint>

[[nodiscard]] std::vector<int> sievePrimes(int max_n) {
    if (max_n < 2) return {};

    // 1. Initialize boolean vector: true means potentially prime
    std::vector<bool> is_prime(max_n + 1, true);
    is_prime[0] = is_prime[1] = false;

    // 2. Mark multiples from p^2 onwards
    for (long long p = 2; p * p <= max_n; ++p) {
        if (is_prime[p]) {
            for (long long multiple = p * p; multiple <= max_n; multiple += p) {
                is_prime[multiple] = false;
            }
        }
    }

    // 3. Collect all primes
    std::vector<int> primes;
    for (int i = 2; i <= max_n; ++i) {
        if (is_prime[i]) primes.push_back(i);
    }
    return primes;
}
\end{lstlisting}

\subsection{What is the Smallest Prime Factor (SPF) Precomputation?}
For applications where you must factorize thousands of numbers up to $N \le 10^7$, we compute an array where \texttt{spf[x]} stores the smallest prime dividing $x$. This enables instant $O(\log x)$ prime factorization.

\begin{lstlisting}[caption={Linear Sieve with Smallest Prime Factor (SPF) for Fast Queries}]
#include <vector>
#include <utility>

class FastFactorizer {
private:
    std::vector<int> spf; // spf[x] = smallest prime factor of x

public:
    explicit FastFactorizer(int max_n) : spf(max_n + 1) {
        // 1. Initialize spf[i] = i
        for (int i = 0; i <= max_n; ++i) spf[i] = i;

        // 2. Mark smallest prime factors
        for (long long i = 2; i * i <= max_n; ++i) {
            if (spf[i] == i) { // i is prime
                for (long long j = i * i; j <= max_n; j += i) {
                    if (spf[j] == j) { // Only set if not already set by a smaller prime
                        spf[j] = i;
                    }
                }
            }
        }
    }

    // O(log x) factorization query
    [[nodiscard]] std::vector<std::pair<int, int>> factorize(int x) const {
        std::vector<std::pair<int, int>> factors;
        while (x > 1) {
            int p = spf[x];
            int count = 0;
            while (x % p == 0) {
                ++count;
                x /= p;
            }
            factors.emplace_back(p, count);
        }
        return factors;
    }
};
\end{lstlisting}

\subsection{What is the Euler Totient Calculation?}
To compute $\phi(n)$ for a single number in $O(\sqrt{n})$ time, apply the formula $\phi(n) = n \prod (1 - 1/p)$. To compute $\phi$ for all numbers up to $N$, adapt the Sieve: initialize $\phi[i] = i$, and whenever a prime $p$ is seen, multiply all its multiples by $(1 - 1/p) = (p - 1)/p$.

\begin{lstlisting}[caption={Single Query $\phi(n)$ and Batch Totient Sieve}]
#include <vector>
#include <cstdint>

// O(sqrt(n)) single query
[[nodiscard]] int64_t computeTotient(int64_t n) {
    int64_t result = n;
    for (int64_t p = 2; p * p <= n; ++p) {
        if (n % p == 0) {
            while (n % p == 0) n /= p;
            result -= result / p; // Equivalent to result * (1 - 1/p)
        }
    }
    if (n > 1) {
        result -= result / n;
    }
    return result;
}

// O(N log log N) batch precomputation for phi(1) ... phi(N)
[[nodiscard]] std::vector<int> totientSieve(int max_n) {
    std::vector<int> phi(max_n + 1);
    for (int i = 0; i <= max_n; ++i) phi[i] = i;

    for (int p = 2; p <= max_n; ++p) {
        if (phi[p] == p) { // p is prime
            for (int multiple = p; multiple <= max_n; multiple += p) {
                phi[multiple] -= phi[multiple] / p;
            }
        }
    }
    return phi;
}
\end{lstlisting}

% ============================================================================
\section{Problem Pattern Recognition}
% ============================================================================

When facing an algorithmic problem, use this step-by-step decision framework to choose the correct number theory tool.

\subsection*{Step 1: What is the Problem Asking For?}
\begin{itemize}
    \item \textbf{``Is this number prime?''} $\implies$ Single query: $O(\sqrt{n})$ trial division with $6k \pm 1$. Range query: Sieve of Eratosthenes.
    \item \textbf{``Find all divisors / count divisors''} $\implies$ Small single number: extract pairs up to $\sqrt{n}$. Counting only: factorize and compute $\prod (a_i + 1)$.
    \item \textbf{``Find the Greatest Common Factor / Simplify Fractions''} $\implies$ Euclidean Algorithm (\texttt{std::gcd}).
    \item \textbf{``Find the smallest number divisible by both $A$ and $B$''} $\implies$ Least Common Multiple (\texttt{std::lcm}).
    \item \textbf{``Count how many numbers $\le N$ share no factors with $N$''} $\implies$ Euler's Totient function $\phi(n)$.
    \item \textbf{``Many queries to factorize numbers''} $\implies$ Precompute SPF array with Sieve, then factorize each in $O(\log n)$.
\end{itemize}

\subsection*{Step 2: What is the Key Constraint?}
\begin{itemize}
    \item \textbf{Query Count vs. Input Size:} Single query with $n \le 10^{12} \implies$ Trial division $O(\sqrt{n})$ runs in $10^6$ steps. Multiple queries ($Q \approx 10^5$) with $n \le 10^7 \implies$ Sieve / SPF precomputation is mandatory.
    \item \textbf{Range $[L, R]$ with large bounds ($R \le 10^{12}$):} If $R - L \le 10^6$, use a Segmented Sieve. Never allocate an array of size $10^{12}$.
    \item \textbf{Extremely large $N \ge 10^{18}$:} Standard trial division fails. Specialized randomized methods (Miller-Rabin for primality, Pollard's Rho for factorization) are required.
\end{itemize}

\subsection*{Quick Reference: Keyword-to-Tool Mapping}

\vspace{0.3cm}
\noindent
\begin{tabularx}{\textwidth}{|p{4.2cm}|X|p{4cm}|}
\hline
\textbf{Keywords in the Problem} & \textbf{Plain English Meaning} & \textbf{Tool to Use} \\ \hline
``is $N$ prime?'' (single query) & Test if any number divides $N$ & Trial division up to $\sqrt{N}$ ($6k \pm 1$) \\ \hline
``count primes up to $N$'' ($N \le 10^7$) & Find all primes in a large range & Sieve of Eratosthenes \\ \hline
``primes in range $[L, R]$'' ($R \le 10^{12}$) & Memory too small for full sieve & Segmented Sieve ($O(\sqrt{R})$ base primes) \\ \hline
``all divisors of $N$'' & Pairs of numbers multiplying to $N$ & $O(\sqrt{N})$ pair collection \\ \hline
``number of divisors of $N$'' & Count total divisors without listing & Prime factorize $\implies \prod (a_i + 1)$ \\ \hline
``greatest common divisor / reduce fraction'' & Largest shared factor between two numbers & Euclidean Algorithm (\texttt{std::gcd}) \\ \hline
``events repeat every $A$ and $B$ days'' & Smallest time when both events align & Least Common Multiple (\texttt{std::lcm}) \\ \hline
``count numbers $\le N$ coprime to $N$'' & Numbers sharing no common prime factors & Euler's Totient $\phi(N) = N \prod (1 - 1/p)$ \\ \hline
``factorize $10^5$ different numbers'' & Individual $O(\sqrt{N})$ tests would TLE & Precompute SPF array, query in $O(\log N)$ \\ \hline
\end{tabularx}
\vspace{0.3cm}

\subsection*{Step 3: Is the Input Too Large for Direct Computation?}
\begin{itemize}
    \item $N \le 10^7$: Array allocation ($10^7$ booleans $\approx 10$ MB) fits in memory. Use a \textbf{standard Sieve}.
    \item $10^7 < N \le 10^{14}$: Array will cause Memory Limit Exceeded. A single query using \textbf{$O(\sqrt{N})$ trial division} takes $\le 10^7$ iterations and runs in $\approx 0.02$ seconds.
    \item $N > 10^{14}$: Trial division reaches $10^8$ operations and risks Time Limit Exceeded. Use Miller-Rabin deterministic primality testing for 64-bit integers.
\end{itemize}

% ============================================================================
\section{Summary and Complexity Reference}
% ============================================================================

\vspace{0.3cm}
\noindent
\begin{tabularx}{\textwidth}{|l|l|l|X|}
\hline
\textbf{Task} & \textbf{Time} & \textbf{Space} & \textbf{Use When} \\ \hline
Primality Test ($6k \pm 1$) & $O(\sqrt{N})$ & $O(1)$ & Single number $N \le 10^{12}$ \\ \hline
All Divisors & $O(\sqrt{N})$ & $O(\sqrt{N})$ & Finding all factors of a single number \\ \hline
Prime Factorization & $O(\sqrt{N})$ & $O(\log N)$ & Single number factorization \\ \hline
$\gcd(a, b)$ / $\operatorname{lcm}(a, b)$ & $O(\log \min(a, b))$ & $O(1)$ & Computing common factors or multiples \\ \hline
Sieve of Eratosthenes & $O(N \log \log N)$ & $O(N)$ & Finding all primes up to $N \le 10^7$ \\ \hline
SPF Precomputation & $O(N \log \log N)$ & $O(N)$ & $Q \ge 10^4$ factorization queries on $N \le 10^7$ \\ \hline
SPF Query Factorization & $O(\log N)$ & $O(\log N)$ & After precomputing SPF array \\ \hline
Euler's Totient $\phi(N)$ & $O(\sqrt{N})$ & $O(1)$ & Single number totient query \\ \hline
Totient Sieve & $O(N \log \log N)$ & $O(N)$ & Computing $\phi(1) \dots \phi(N)$ in batch \\ \hline
Segmented Sieve & $O((R - L) \log \log R)$ & $O(R - L)$ & Primes in $[L, R]$ with $R \le 10^{12}$, $R-L \le 10^6$ \\ \hline
\end{tabularx}
\vspace{0.3cm}

% ============================================================================
\section{Glossary of Terms and Symbols}
% ============================================================================

If you are new to number theory, this dictionary provides plain-English definitions for every term and symbol used in this document:

\subsection*{Core Concepts}
\begin{itemize}
    \item \textbf{Prime Number}: A whole number greater than 1 that cannot be formed by multiplying two smaller whole numbers.
    \item \textbf{Composite Number}: A whole number that has divisors other than 1 and itself (for example, $6 = 2 \times 3$).
    \item \textbf{Factor / Divisor}: A whole number that divides another evenly without leaving a remainder.
    \item \textbf{Coprime (Relatively Prime)}: Two numbers that share no common factors other than 1 (for example, 8 and 15).
    \item \textbf{Prime Factorization}: Writing a number as a product of prime powers (for example, $12 = 2^2 \times 3$).
    \item \textbf{Greatest Common Divisor (GCD)}: The largest positive integer that divides both numbers evenly.
    \item \textbf{Least Common Multiple (LCM)}: The smallest positive integer that is a multiple of both numbers.
    \item \textbf{Trial Division}: The direct method of testing whether a number $n$ is prime by dividing it by every candidate factor from 2 up to $\sqrt{n}$. The $6k \pm 1$ optimization skips all multiples of 2 and 3, reducing the number of tests by two thirds.
    \item \textbf{Sieve of Eratosthenes}: An algorithm that finds all primes up to $N$ by systematically crossing out multiples. Much faster than calling trial division on every number individually.
    \item \textbf{Segmented Sieve}: A variant of the Sieve of Eratosthenes that works on a range $[L, R]$ without allocating memory proportional to $R$. Used when $R$ is too large for a standard sieve array but $R - L$ is manageable.
    \item \textbf{Smallest Prime Factor (SPF)}: The smallest prime number that divides $x$. Used to factorize numbers in $O(\log x)$.
    \item \textbf{Euler's Totient Function $\phi(n)$}: Counts how many numbers from 1 to $n$ share no common factors with $n$.
    \item \textbf{Multiplicative Function}: A number-theoretic function where $f(a \cdot b) = f(a) \cdot f(b)$ whenever $\gcd(a, b) = 1$. Both $\phi(n)$, $\tau(n)$, and $\sigma(n)$ are multiplicative.
\end{itemize}

\subsection*{Mathematical Symbols}
\begin{itemize}
    \item $a \mid b$ \textbf{(Divides)}: Means ``$a$ divides $b$ with zero remainder'' (for example, $3 \mid 12$).
    \item $a \nmid b$ \textbf{(Does Not Divide)}: Means ``$a$ does not divide $b$ evenly'' (for example, $5 \nmid 12$).
    \item $\gcd(a, b)$: The Greatest Common Divisor of $a$ and $b$.
    \item $\operatorname{lcm}(a, b)$: The Least Common Multiple of $a$ and $b$.
    \item $\phi(n)$ \textbf{(Euler's Totient)}: The count of integers in $[1, n]$ coprime to $n$.
    \item $\tau(n)$ or $d(n)$ \textbf{(Number of Divisors)}: The total count of positive divisors of $n$.
    \item $\sigma(n)$ \textbf{(Sum of Divisors)}: The sum of all positive divisors of $n$.
    \item $\prod$ \textbf{(Product Notation)}: The Greek letter capital Pi. Means ``multiply all these terms together in a loop'' (just as $\sum$ means add).
    \item $\lfloor x \rfloor$ \textbf{(Floor)}: Rounding a decimal down to the nearest whole integer (for example, $\lfloor 3.8 \rfloor = 3$).
    \item $\sqrt{n}$ \textbf{(Square Root)}: The number that, when multiplied by itself, equals $n$.
\end{itemize}

% ============================================================================
\section{Further Reading}
% ============================================================================

\subsection*{Free Online Resources (Start Here)}
\begin{itemize}
    \item \textbf{CP-Algorithms: Number Theory} (\url{https://cp-algorithms.com}) --- High-performance implementations of the Sieve, Linear Sieve, GCD, and Euler's Totient function with rigorous algorithmic complexity notes.
    \item \textbf{Art of Problem Solving (AoPS) Number Theory Wiki} (\url{https://artofproblemsolving.com/wiki}) --- Clear, intuitive explanations of divisibility, prime factorization, and modular arithmetic from competitive math trainers.
    \item \textbf{Brilliant.org: Number Theory} (\url{https://brilliant.org}) --- Interactive, visual puzzles covering primes, factorization, and modular arithmetic from intuition to mastery.
    \item \textbf{Project Euler} (\url{https://projecteuler.net}) --- The premier training ground for number theory. Highly recommended starter problems directly solved by this guide:
    \begin{itemize}
        \item \textbf{Problem 3}: Largest Prime Factor (Trial division / factorization).
        \item \textbf{Problem 5}: Smallest Multiple (LCM across an array).
        \item \textbf{Problem 7}: 10001st Prime (Sieve of Eratosthenes).
        \item \textbf{Problem 10}: Summation of Primes (Sieve summation).
        \item \textbf{Problem 12}: Highly Divisible Triangular Number (Divisor counting $\tau(n)$ via prime factorization).
        \item \textbf{Problem 69}: Totient Maximum (Euler's totient function properties).
        \item \textbf{Problem 72}: Counting Fractions (Sum of Euler's totient values).
    \end{itemize}
\end{itemize}

\subsection*{Books (When You Are Ready to Go Deeper)}
\begin{itemize}
    \item K.\ H.\ Rosen, \emph{Elementary Number Theory and Its Applications} --- an approachable, well-written undergraduate text covering prime numbers, divisibility, the Euclidean algorithm, and cryptography from scratch.
    \item G.\ H.\ Hardy and E.\ M.\ Wright, \emph{An Introduction to the Theory of Numbers} --- the timeless classic of mathematical number theory. Dense, elegant, and exhaustive in its coverage of primes and divisors. Read this after Rosen, not before.
    \item S.\ Halim and F.\ Halim, \emph{Competitive Programming} --- dedicated chapter on Mathematics covers sieves, GCD, and number-theoretic algorithms with contest-tested C++ snippets.
\end{itemize}

% --- Version History (always the last section of the document) ---
\section*{Version History}
\addcontentsline{toc}{section}{Version History}

\noindent
\begin{tabularx}{\textwidth}{|l|X|}
\hline
\textbf{Date} & \textbf{Change} \\ \hline
2026-09-05 & Document created. \\ \hline
2026-10-03 & Refactored against the learning materials skill: canonical preamble, worked examples added to GCD, LCM, divisor-count and divisor-sum subsections, traces merged into Key Algorithmic Patterns, intuition callouts and text diagrams added, banned phrases and platform names removed. \\ \hline
2026-10-03 & Fixed integer overflow bugs in Sieve of Eratosthenes and FastFactorizer loops. \\ \hline
\end{tabularx}

\end{document}
